\documentclass[9pt]{article}

\usepackage[margin=0.2in]{geometry}
\usepackage[T1]{fontenc}
\usepackage[utf8]{inputenc}
\usepackage{lmodern}
\usepackage{listings}
\usepackage{xcolor}
\usepackage{hyperref}

\setlength{\parindent}{0pt}  % <-- removes all paragraph indentation

% Simple JavaScript definition for listings (enough for our examples)
\lstdefinelanguage{JavaScript}{
  keywords={break, case, catch, continue, debugger, default, delete, do, else,
    finally, for, function, if, in, instanceof, new, return, switch, this, throw,
    try, typeof, var, void, while, with, const, let},
  sensitive=true,
  morecomment=[l]{//},
  morecomment=[s]{/*}{*/},
  morestring=[b]",
  morestring=[b]'
}

% Highlight 'fields' and 'allFields' differently inside code
\lstset{
  basicstyle=\ttfamily\small,
  breaklines=true,
  columns=fullflexible,
  showstringspaces=false,        % hides visible space markers in strings
  emph={fields},                      % group 1
  emphstyle=\color{blue}\bfseries,
  emph={[2]allFields},                % group 2
  emphstyle={[2]\color{red}\bfseries}
}

\begin{document}

\noindent
\textbf{Database Management Systems}%
\hfill%
\textit{Creating a MongoDB Replica Set on a Single Machine}\\
Beata Kubiak, Instructor

\medskip

A MongoDB ``cluster'' is made of \texttt{mongod} processes (servers), not
necessarily physical machines. Multiple MongoDB servers can run on the same
computer, as long as each process uses its own data directory.

Each MongoDB server process (each \texttt{mongod}) needs its own separate data
directory because:
\begin{itemize}
    \item They are separate servers, even when running on the same machine.
    \item Each server stores its own physical copy of the replicated data.
    \item If they shared the same directory, they would corrupt each other's files.
\end{itemize}

This setup allows us to fully simulate a multi-node cluster on a single machine.

\section*{1. Create Data Directories}

Create three directories for the three replica-set members:

\begin{verbatim}
e:\mongodb\rs1\
e:\mongodb\rs2\
e:\mongodb\rs3\
\end{verbatim}

Each directory belongs to one MongoDB node.

\section*{2. Start Three mongod Instances}

Start each server in its own terminal:

\begin{verbatim}
mongod.exe --port 27017 --dbpath e:\mongodb\rs1 --replSet rs0
mongod.exe --port 27018 --dbpath e:\mongodb\rs2 --replSet rs0
mongod.exe --port 27019 --dbpath e:\mongodb\rs3 --replSet rs0
\end{verbatim}

All three processes share the replica set name \texttt{rs0}, which tells MongoDB
they should eventually join the same replica set.

\section*{3. Initiate the Replica Set}

Connect to one of the nodes (for example, port 27017):

\begin{verbatim}
mongosh --port 27017
\end{verbatim}

Then run:

\begin{verbatim}
rs.initiate({
  _id: "rs0",
  members: [
    { _id: 0, host: "localhost:27017" },
    { _id: 1, host: "localhost:27018" },
    { _id: 2, host: "localhost:27019" }
  ]
})
\end{verbatim}

This command tells MongoDB the exact list of servers that belong to the replica set,
establishes communication between the nodes, and triggers an election of a primary
and secondaries.

\section*{4. Replica Set Behavior}

Once the replica set is initiated:
\begin{itemize}
    \item The nodes exchange heartbeats.
    \item They share their oplogs.
    \item They elect a primary.
    \item The remaining nodes become secondaries.
    \item The nodes synchronize data.
    \item Read/write concerns begin to apply.
    \item Failover becomes possible.
\end{itemize}

At this point, you have a three-node replica set running on one machine.

\section*{5. Restoring Data for Replication}

To make all three nodes contain the same database, restore the data only on the
primary node (typically the one on port 27017). The secondaries automatically
replicate the data from the primary.

To restore only the \texttt{sample\_mflix} database from your archive:

\begin{verbatim}
mongorestore --port=27017 --archive=sampledata.archive --nsInclude="sample_mflix.*"
\end{verbatim}

This writes the \texttt{sample\_mflix} data into the primary. After the restore
completes, the secondaries detect the changes via the oplog and copy the data
automatically.


\end{document}
